% Problema: comprobar axiomas de norma (castellano).
%
% Enunciado, respuesta, solución y corrección en UN fichero. Cada nivel
% aparece solo en el perfil que corresponde: el mismo fichero produce la hoja
% del alumno, la hoja con respuestas, la hoja con soluciones y la del profesor.

\begin{exercise}[Axiomas de norma]
Decide cuáles de las siguientes funciones $\mathbb{R}^2 \to \mathbb{R}$ son
normas. Para las que no lo sean, indica qué axioma falla.

\begin{parts}
  \item $f(x_1, x_2) = |x_1| + 2|x_2|$
  \item $f(x_1, x_2) = |x_1| - |x_2|$
  \item $f(x_1, x_2) = \sqrt{x_1^2 + x_2^2}$
  \item $f(x_1, x_2) = \max\{|x_1|, |x_2|\}^2$
\end{parts}

\dmarks{4}

\begin{hint}
Para descartar una candidata basta encontrar un contraejemplo de un solo
axioma. Prueba con vectores sencillos y con $\lambda = 2$.
\end{hint}

\begin{answer}
(i) sí; (ii) no, falla la positividad; (iii) sí; (iv) no, falla la homogeneidad.
\end{answer}

\begin{solution}
\begin{parts}
  \item Es una norma. Es la norma 1 con pesos: los tres axiomas se comprueban
        término a término, y la desigualdad triangular se sigue de la de
        $|\cdot|$ en cada coordenada.
  \item No es una norma. Falla la positividad: $f(1,1) = 0$ con $(1,1) \neq 0$.
        También falla la desigualdad triangular.
  \item Es la norma euclídea.
  \item No es una norma. Falla la homogeneidad:
        \[
          f(2x) = \max\{2|x_1|, 2|x_2|\}^2 = 4\max\{|x_1|,|x_2|\}^2 = 4 f(x),
        \]
        mientras que la homogeneidad exigiría $f(2x) = 2 f(x)$.
\end{parts}
\end{solution}

\begin{marking}
Un punto por apartado. En (ii) y (iv) se exige nombrar el axioma que falla; dar
solo "no es norma" vale medio punto. El error habitual en (iv) es comprobar la
desigualdad triangular y no la homogeneidad.
\end{marking}
\end{exercise}

\begin{exercise}[Normas equivalentes]
Demuestra que en $\mathbb{R}^n$ se cumple
\[
  \|x\|_\infty \leq \|x\|_2 \leq \sqrt{n}\,\|x\|_\infty .
\]
\dmarks{3}

\begin{answer}
Ambas desigualdades salen de acotar cada $|x_i|$ por el máximo.
\end{answer}

\begin{solution}
Sea $m = \|x\|_\infty = \max_i |x_i|$ y sea $k$ un índice donde se alcanza.
Para la primera desigualdad,
\[
  \|x\|_2^2 = \sum_{i=1}^n |x_i|^2 \geq |x_k|^2 = m^2 ,
\]
y tomando raíces, $\|x\|_2 \geq m$. Para la segunda, $|x_i| \leq m$ para todo
$i$, luego
\[
  \|x\|_2^2 = \sum_{i=1}^n |x_i|^2 \leq n m^2 ,
\]
y por tanto $\|x\|_2 \leq \sqrt{n}\, m$.
\end{solution}
\end{exercise}
